\chapter{Introdu��o}
\label{intro}

Escrever a introdu��o neste arquivo.

O Cap�tulo \ref{intro} � a ``Introdu��o''.

O Cap�tulo \ref{del} refere-se a ``Delimita��o do Tema''.

O Cap�tulo \ref{prob} refere-se a ``Problemas e Premissas''.


A Figura \ref{fig:elem1} ocupa toda a largura do texto.

\begin{figure}[!htbp]
	\centering
	\captionsetup{width=0.50\textwidth} % Tamanho da legenda da figura (70% da largura do texto).
		\includegraphics[width=1.00\textwidth]{../figuras/elementos.eps}
	\caption{Primeira Figura.}
		\label{fig:elem1}
\end{figure}

A Figura \ref{fig:elem2} ocupa metade da largura do texto.

\begin{figure}[!htbp]
	\centering
		\captionsetup{width=0.50\textwidth} % Tamanho da legenda da figura (70% da largura do texto).
		\includegraphics[width=.50\textwidth]{../figuras/elementos.eps}
	\caption{Segunda Figura.}
		\label{fig:elem2}

\end{figure}

A Figura \ref{fig:elem3} ocupa um quarto da largura do texto.
\begin{figure}[!htbp]
	\centering
	\captionsetup{width=0.50\textwidth} % Tamanho da legenda da figura (70% da largura do texto).
		\includegraphics[width=.25\textwidth]{../figuras/elementos.eps}
	\caption{Terceira Figura}
		\label{fig:elem3}
\end{figure}

A Figura \ref{fig:elem3} ocupa um oitavo da largura do texto.
\begin{figure}[!htbp]
	\centering
	\captionsetup{width=0.50\textwidth} % Tamanho da legenda da figura (70% da largura do texto).
		\includegraphics[width=.125\textwidth]{../figuras/elementos.eps}
	\caption{Quarta Figura}
		\label{fig:elem4}
\end{figure}


\clearpage

A Tabela \ref{tab:t1}  � apresentada abaixo.
\begin{table}[!htbp]
	\centering
	\captionsetup{width=0.30\textwidth} % Tamanho da legenda da figura (30% da largura do texto).
	\caption{Primeira tabela}
\begin{tabular}{llllll}
\hline
\multicolumn{1}{|l|}{e} & \multicolumn{1}{l|}{a1} & \multicolumn{1}{l|}{a2} & \multicolumn{1}{c|}{teste} & \multicolumn{1}{r|}{f} & \multicolumn{1}{r|}{z} \\ 
\hline
\multicolumn{1}{|l|}{12} & \multicolumn{1}{l|}{} & \multicolumn{1}{l|}{} & \multicolumn{1}{l|}{} & \multicolumn{1}{r|}{} & \multicolumn{1}{r|}{} \\ 
\hline
3 &  & 5 & r &  &  \\ 
4 & 5\% &  &  &  &  \\ 
\hline
\end{tabular}

	\label{tab:t1}
\end{table}


A Tabela \ref{tab:t2} � apresentada abaixo.
\begin{table}[!htbp]
	\centering
	\captionsetup{width=0.3\textwidth} % Tamanho da legenda da figura (30% da largura do texto).
	\caption{Segunda tabela}
\begin{tabular}{|l|l|l|l|}
\hline
\multicolumn{1}{|c|}{} & \multicolumn{1}{c|}{Acelera��o (m/s)} & \multicolumn{1}{c|}{Tempo (s)} & \multicolumn{1}{c|}{} \\ 
\hline
\multicolumn{1}{|c|}{} & \multicolumn{1}{c|}{0} & \multicolumn{1}{c|}{0} & \multicolumn{1}{c|}{} \\ 
\hline
\multicolumn{1}{|c|}{} & \multicolumn{1}{c|}{1} & \multicolumn{1}{c|}{10} & \multicolumn{1}{c|}{} \\ 
\hline
\multicolumn{1}{|c|}{} & \multicolumn{1}{c|}{2} & \multicolumn{1}{c|}{20} & \multicolumn{1}{c|}{} \\ 
\hline
\end{tabular}

	\label{tab:t2}
\end{table}


sendo que \simb[Di�metro da circunfer�ncia]{\alpha} � o di�metro da circunfer�ncia.

\noindent Este s�mbolo (\simb[Gradiente de uma fun��o]{\nabla}) foi inserido no texto, fora do ambiente \textit{equation}.



Exemplos de siglas:
A organiza��o das na��es unidas \sigla[Organiza��o das Na��es Unicas]{ONU};
\sigla[Teste da lista de siglas 2. Esta possui um significado bem grande que ocupa mais de uma linha]{EASREFER};
\sigla[Teste da lista de siglas 3]{ER}.

A primeira refer�ncia \cite{abntex-doc} � um livro. A segunda refer�ncia \cite{anil2004} � um artigo e peri�dico. A terceira � um link \cite{matlab}. 
